% Chapter 4

\chapter{Introdução} % Main chapter title
\label{chap:Chapter4} % For referencing the chapter elsewhere, use \ref{chap:Chapter4} 
Este capítulo apresenta a organização responsável pela ideia do tema a ser desenvolvido, o contexto do seu surgimento e os problemas que a originaram. São abordados os objetivos, o planeamento e as considerações éticas tidas em conta para executar a proposta. Por último é descrita a estrutura deste documento de forma a facilitar a leitura do mesmo.
%----------------------------------------------------------------------------------------
\section{Organização e Contexto}
A Critical Techworks (CTW) nasceu de um empreendimento conjunto entre a Critical Software e o grupo BMW (Bayerische Motoren Werke). Estabelecida em 2018 esta empresa desenvolve software exclusivamente para os produtos da BMW. A sua missão é a de "transformar a forma como o mundo se move" focando-se para isso na transformação digital com as melhores práticas do mercado.

A organização estrutura-se em equipas ágeis multidisciplinares, agrupadas em unidades e domínios tecnológicos, o que promove autonomia, ciclos de desenvolvimento rápidos e colaboração contínua. Esta abordagem permite alinhar o desenvolvimento de software com necessidades reais dos veículos e serviços digitais da BMW.

O presente relatório foi desenvolvido no âmbito da unidade curricular de Preparação para Dissertação (PREPD) do Instituto Superior de Engenharia do Porto (ISEP), com o intuito de demonstrar o estudo de viabilidade de uma solução para o problema em causa.

\section{Problema}

Na área de TI (Tecnologias de Informação), as equipas podem manter vários serviços em simultâneo e por essa razão os engenheiros de plantão (membros das equipas responsáveis por manter os serviços estáveis \textbf{24x7}) podem ficar \textbf{sobrecarregados} de incidentes, fazendo com que tenham que lidar com eventos provenientes de vários serviços paralelamente. Numa situação como esta, os engenheiros podem ter um sentimento de \textbf{confusão}, misturando temas e errando nas suas análises, o que pode causar piores resoluções de incidentes e o aumento do seu tempo de mitigação (\textit{\textbf{TTM}}). Outro cenário frequente para os \textit{\textbf{OCEs}} (\textit{On-Call Engineers}) envolve incidentes que ocorrem fora do horário de expediente, especialmente durante a madrugada, quando o engenheiro de prevenção pode não estar na sua plena capacidade de raciocínio e operacional.

Foi também referido anteriormente os incidente com origem manual. Estes, são criados por clientes finais e/ou equipas de testes e podem conter mais ou menos informação. Muitas vezes, os \textit{OCEs} são \textbf{obrigados} a fazer uma \textbf{extensa análise} para compreender o conteúdo do incidente e a resolução mais adaptada a aplicar.

Uma interpretação ou resolução inadequada das ocorrências pode \textbf{reduzir a disponibilidade} dos serviços, o que, por sua vez, pode gerar \textbf{prejuízos significativos} para a organização, seja pela perda de receitas, seja pelo impacto negativo na experiência do utilizador final.

\section{Objetivos}
	O objetivo desta investigação passa por perceber que soluções existem atualmente no mercado para solucionar o problema mencionado anteriormente, entender se as soluções existentes são eficazes nos seus contextos e se  \hl{preenchem todos os requisitos necessários para satisfazer todas as partes interessadas}.
	
	- Verificar soluções no mercado
	- Perceber viabilidade de uma solução nova
	- Desenvolvimento de um planeamento para o desenvolvimento da solução
	- Explicar objetivos da solução:
		- Geração de resumos dos incidentes (com base no histórico)
		- Geração de sugestões de mitigação (com base no histórico)
		- (Opcional ^^ Com base em logs)
	- Beneficios esperados (diminiução de carga, ttm...)

\section{Planeamento}
Gant chart...

\section{Considerações Éticas}
	Neste projeto, foram seguidos os princípios dos códigos de ética da \textit{ACM} (\textit{Code of Ethics and Professional Conduct}), \textit{NSPE} (\textit{Code of Ethics for Engineers}) e da \textit{ACM/IEEE-CS} (\textit{Software Engineering Code}), de forma a orientar o estudo a atuar com responsabilidade, tanto em relação aos dados como às pessoas afetadas pela solução.

	Tendo em conta os dados a serem utilizados poderem \textbf{incluir informações confidenciais}, foi avaliada cuidadosamente a possibilidade de anonimizar os dados e, caso não seja satisfatório para as partes interessadas, será garantido que apenas os resultados agregados sejam apresentados, \textbf{evitando quaisquer riscos de exposição de dados confidenciais}. Não identificámos outros riscos éticos relevantes no projeto.
	
	Para além dos códigos acima mencionados, foi também seguido o Código de Boas Práticas e de Conduta do Instituto Politécnico do Porto.
	
	Quanto às competências necessárias, o autor é também um dos muitos engenheiros que lidam com o tipo de problema que este projeto aborda, estando também a aprofundar os conceitos no mestrado que originou esta investigação. Essa combinação de experiência na área e aprendizagem académica ajuda a assegurar que o trabalho é feito com rigor, ética e atenção às melhores práticas profissionais.
	
	\hl{Refinar}

\section{Estrutura do Documento}
